\section{Introduction}
\label{sec:introduction}

Standard feedforward neural networks organize computation along a linear sequence of layers, and their training is classically formulated through backpropagation \cite{RumelhartHintonWilliams1986}. Modern deep architectures often enrich this linear organization by introducing skip connections, residual pathways, or other mechanisms that improve information propagation across depth \cite{HeZhangRenSun2016}. At the same time, memory-oriented models suggest that useful architectures should support both a stable, slowly integrated pathway and a rapidly accessible pathway for short-lived information \cite{HochreiterSchmidhuber1997,WestonChopraBordes2014,GravesWayneDanihelka2014}.

A different but complementary line of work has emphasized that graphs alone encode only pairwise interactions, whereas simplicial and more general topological structures provide a natural framework for higher-order interactions \cite{BarbarossaSardellitti2019,EbliDefferrardSpreemann2020,BodnarFrascaWangEtAl2021,HajijEtAl2022}. Most of this literature focuses on extending signal processing or message passing to simplicial domains. In the present note, our goal is different: we propose a \emph{simplex-based feedforward memory architecture} whose geometry itself organizes input, output, hidden propagation, and memory coupling.

The basic idea is simple. We begin with a regular $n$-simplex, discretize it by a lattice of points, and view these lattice points as hidden nodes. Two opposite facets of the simplex act as input and output interfaces, while their common $(n-2)$-dimensional face serves as a shared intermediate buffer. A linear potential then induces a directed nearest-neighbor graph on the lattice, yielding a feedforward hidden architecture. In the minimal motifs, the triangle and the tetrahedron already exhibit the qualitative features we seek: a narrow backbone, parallel intermediate channels, and a natural distinction between deep and shallow propagation paths.

The contribution of this note is conceptual and mathematical. First, we define a family of directed hidden graphs $\SMN(n,m)$ based on the lattice of a regular $n$-simplex with $m$ points on each edge. Second, we identify the input facet, output facet, and shared interface, and compute several basic combinatorial quantities associated with these sets. Third, we describe the triangle and tetrahedron motifs explicitly as minimal geometric instances of the construction. Finally, we propose a geometric interpretation in which a narrow deep backbone represents long-term memory, while a broad shallow interface represents short-term memory.
